% !PRINTED none
% Abschnitt-opening page: no running-head folio number is printed on this
% leaf (suppressed on the chapter-opening page, as is common in this book);
% resumes as "564" on the next leaf.

\begin{center}
\gesperrt{Twenty-Third Chapter.}

{\Large \textbf{The Division Field.}}
\end{center}

\begin{center}
\rule{2cm}{0.4pt}
\end{center}

\begin{center}
\S\,150. \textbf{The Homogeneous \gesperrt{Weierstrass} Functions.}
\end{center}

In \S\,114 we set out from the question under what suppositions a doubly
periodic function $\varphi(u)$ with the periods $\omega_1, \omega_2$
admits a multiplication $\varphi(\mu u)$, that is, under what
circumstances $\mu\omega_1, \mu\omega_2$ can be expressed linearly and
with integral coefficients by $\omega_1, \omega_2$; and we had arrived at
the result that, for variable $\omega_1, \omega_2$, this is possible only
when $\mu$ is an integral rational number, and that only when

\weq{1}{\omega = \frac{\omega_2}{\omega_1}}

is an imaginary quadratic irrational number can $\mu$ also be a complex
number of the same field as $\omega$. What now matters is to explore
somewhat more closely the formulae of complex multiplication that follow
from this assumption.

We will first investigate, as the formally simpler case, the complex
multiplication of the \gesperrt{Weierstrass} elliptic function $\wp(u)$, but we must
then also take into consideration the complex multiplication of the
\gesperrt{Jacobi} elliptic functions.

In \S\,37, (8), \S\,46, (13) we derived the formulae for the
homogeneity of the functions $\sigma(u)$, $\wp(u)$ and of the invariants
$g_2, g_3, 16\,G = g_2^3 - 27\,g_3^2$, which we here collect once more:

\weq{2}{\begin{aligned}
\sigma(\lambda u, \lambda\omega_1, \lambda\omega_2) &= \lambda\ \sigma(u, \omega_1, \omega_2), \\
\wp(\lambda u, \lambda\omega_1, \lambda\omega_2) &= \lambda^{-2}\, \wp(u, \omega_1, \omega_2), \\
g_2(\lambda\omega_1, \lambda\omega_2) &= \lambda^{-4}\, g_2(\omega_1, \omega_2), \\
g_3(\lambda\omega_1, \lambda\omega_2) &= \lambda^{-6}\, g_3(\omega_1, \omega_2), \\
G(\lambda\omega_1, \lambda\omega_2) &= \lambda^{-12}\, G(\omega_1, \omega_2),
\end{aligned}}
%CONT>>
